\section{Conclusion}
\label{sec:conclusion}

We have argued that the architecturally consequential dimension of forensic
evidence is not what the evidence is, but how it is addressed and
traversed. We presented a taxonomy of five navigation primitives --
persistent storage \src{P}, memory \src{M}, log \src{L}, live query
\src{Q}, and content-addressed \src{C} -- and developed its formal
consequences. We proved that parsers operating on byte sequences compose
uniformly with all five primitives (Theorem~\ref{thm:parser-indep}); we
derived from that theorem a layer-dependency rule that constrains how
forensic software is to be organized (Theorem~\ref{thm:layer-rule}); we
formalized attacker-durability and characterized its distinct profiles
across the primitives (Definition~\ref{def:durable},
Propositions~\ref{prop:q-durability}~and~\ref{prop:c-durability},
Corollary~\ref{cor:c-anchor-free}); and we presented Issen, a Rust
reference implementation in which the layer rules are enforced at the
crate boundary.

The taxonomy provides the formal basis for a kind of cross-source
correlation that has been an open challenge since Garfinkel's 2010
agenda~\cite{garfinkel2010digital}. By distinguishing source types by
navigation primitive rather than by artifact category, a correlation
engine can route evidence to a parser-of-record that is independent of
the source (Theorem~\ref{thm:parser-indep}), can apply source-appropriate
trust weights (Propositions~\ref{prop:q-durability}~and~\ref{prop:c-durability},
Corollary~\ref{cor:c-anchor-free}), and can pivot across sources -- in
particular, can pivot across the cryptographic-hash address space of
\src{C} (Definition~\ref{def:hash-pivot}) -- in a principled way.

We have made specific architectural commitments: parsers must accept
byte sequences and must not depend on navigation crates; live-query results
must be transmitted off-host promptly to acquire durability; content-addressed
artifacts can be trusted by hash equality without external chain-of-custody
infrastructure. Each commitment follows from the formal properties; each
has been instantiated in deployed code; each carries forensic
consequences that artifact-type taxonomies do not capture.

Several extensions remain. A formal completeness theorem for the five
primitives is open. The \src{C} implementation in Issen is in development.
Empirical evaluation across NIST and DFRWS standard datasets is planned.
We expect future work in cross-source correlation algorithms and in
adversarial-model parameterization to build on the foundations laid here.

The premise of digital forensics has always been that evidence persists
even when actors prefer it not to. The five primitives are five distinct
modes by which that persistence is realized: by writing to disk, by
residing in memory, by accumulating in logs, by being captured into
investigator hands, by being identified with an unforgeable hash. Each
mode has its own properties and its own discipline. The taxonomy is the
field's formal acknowledgement of that plurality.
